\documentclass[10pt]{article}

\usepackage[utf8]{inputenc}

\usepackage[T1]{fontenc}

\usepackage{amsmath}

\usepackage{amsfonts}

\usepackage{amssymb}

\usepackage[version=4]{mhchem}

\usepackage{stmaryrd}

\usepackage{graphicx}

\usepackage[export]{adjustbox}

\graphicspath{ {./images/} }

\usepackage{bbold}



\begin{document}

(см. Пенлеве уравнения); все они гамильтоновы и имеют коммутационное представление типа лаксовой пары.



Важную роль играют асимптотич. представления в окрестности существенно особых точек решений этих шести канонич. типов уравнений Пенлеве в терминах их первых интегралов, в качестве к-рых выступают данные монодромии вспомогательной линейной системы, коэффициентами к-рой являются функции Пенлеве и их производные. Решение прямой и обратной задачи теории монодромии для такой линейной системы эквивалентно построению интегрального представления функции Пенлеве. Данные монодромии вычисляются явным образом по параметрам асимптотик, что позволяет получить явные формулы связи асимптотик в различных существенно особых точках или в различных областях вблизи одной особой точки.



Лuт.: [1] А й н с Э. Л., Обыкновенные дифференциальные уравнения, пер. с англ., Харьков, 1939; [2] It s A. R., Novokshe n ov V. Yu., «Lecture Notes Math.», v. 1191, B., 1986.



В. Ю. Новокшенов. ПEHЛЕВЕ УPABНЕНИЯ - общее название группы из шести обыкновенных дифференциальных уравнений. Введены П. Пенлеве (P. Painlevé, 1900) и Б. Гамбье (B. Gambier, 1910) при классификации уравнений типа



\[

w^{\prime \prime}=R\left(z, w, w^{\prime}\right),

\]



где $R$ - функция аналитическая по $z$ и рациональная по $w$ и $w^{\prime}$.



Обычно П. у. записывают в следующем виде:



\[

\begin{aligned}

\text { I. } w^{\prime \prime} & =6 w^{2}+z \\

\text { II. } w^{\prime \prime} & =2 w^{3}+z w+a \\

\text { III. } w^{\prime \prime} & =\left(w^{\prime}\right)^{2} w^{-1}+\left(a w^{2}+b\right) \exp z+ \\

+ & \left(c w^{3}+d w^{-1}\right) \exp 2 z,|b|+|d| \neq 0, \\

\text { IV. } w^{\prime \prime} & =\left(w^{\prime}\right)^{2} w^{-1} / 2+3 w^{3} / 2+4 z w^{2}+2\left(z^{2}-\alpha\right) w+\beta w^{-1}, \\

\text { V. } w^{\prime \prime} & =\left(w^{\prime}\right)^{2}\left(\frac{1}{2 w}+\frac{1}{w-1}\right)-\frac{w^{\prime}}{z}+ \\

+ & \frac{(w-1)^{2}}{z^{2}}\left(\alpha w+\frac{\beta}{w}\right)+\gamma \frac{w}{z}+\delta \frac{w(w+1)}{w-1}, \\

\text { VI. } w^{\prime \prime} & =\frac{\left(w^{\prime}\right)^{2}}{2}\left(\frac{1}{w}+\frac{1}{w-1}+\frac{1}{w-z}\right)-w^{\prime}\left(\frac{1}{z}+\frac{1}{z-1}+\frac{1}{w-z}\right)+ \\

+ & \frac{w(w-1)(w-z)}{z^{2}(z-1)^{2}}\left[\alpha+\beta \frac{z}{w^{2}}+\gamma \frac{z-1}{(w-1)^{2}}+\delta \frac{z(z-1)}{(w-z)^{2}}\right] .

\end{aligned}

\]



П. у. возникают при сведении к обыкновенным дифференциальным уравнениям некоторых нелинейных уравнений математич. физики, в частности Кортевега - де Фриса уравнения (П.у. II), синус-Гордона уравнения (П.у. III), нелинейного уравнения ШІредингера (П.у. IV).



Решения П.у. (трансцендентные функции Пенле в е - специальные функции, не сводящиеся к известным) обладают Пенлеве свойством: не имеют других подвижных (то есть зависящих от постоянных интегрирования или начальных данных) особенностей, кроме полюсов. Так, решения П.у. I-IV не имеют вообще никаких особенностей, кроме полюсов; решения П. у. V имеют неподвижные логарифмич. точки ветвления при $z=0$ и $z=\infty$, а решения П.у. VI - при $z=0, z=1$ и $z=\infty$. Установление свойства Пенлеве позволяет находить интегрируемые варианты различных моделей нелинейных явлений и многих нелинейных уравнений, решаемых при помощи обратной задачи рассеяния метода.



Лит.: [1] А й н с Э. Л., Обыкновенные дифференциальные уравнения, пер. с англ., Харьков, 1939; [2] Гол убе в В. В., Лекшии по аналитической теории дифференциальных уравнений, 2 изд., М.- Л., 1950; [3] Арнольд В.И., Ильяшенко Ю. С., в кн.: Итоги науки и техники. Современные проблемы математики, т. 1, М., 1985. Ю. А. Данилов.



ПЕНРОУЗА ДИАГРАММЫ - координатные диаграммы конформного пространства-времени, предложенные Р. Пенроузом [1]. П.д. строятся в системе координат, переход к к-рым делает возможным представление бесконечности в виде границы конечной области. Конформное преобразование пространства-времени позволяет ввести регулярную метрику на расширенном многообразии, образованном присоединением этой границы к исходному многообразию. На рис. изображена П. д. $r=t$ плоскости аналитически про-



\begin{center}

\includegraphics[max width=\textwidth]{2023_07_08_f58722d513d98683e1d1g-1}

\end{center}



долженного пространства-времени Шварцшильда в координатах $\psi, \xi$, связанных с координатами Крускала - Шекерса $u, v$ по формулам



\[

v+u=\operatorname{tg}(\psi+\xi) / 2, v-u=\operatorname{tg}(\psi-\xi) / 2 .

\]



Области I, III $(r>2 m)$ на диаграмме соответствуют двум асимптотически плоским областям, области II, IV $(r<2 m)$ - внутренности черной и белой дыры; горизонты изображаются линиями $\psi= \pm \xi$. Диаграмма содержит две сингулярности $r=0$ и по две изотропных конформных бесконечности $J^{+}$и $J^{-}$(см. Конформная бесконечность асимптотически плоского пространства-времени). Времениподобные бесконечности, где начинаются и кончаются времениподобные кривые и, в частности, мировые линии $r=$ const (изображены на рис. пунктиром), обозначаются на диаграмме точками $i^{-}$и $i^{+}$. Точки $i^{0}$ представляют собой пространственноподобные бесконечности, куда простираются сечения $t=$ const (изображены сплошными линиями). Радиальные нулевые геодезические имеют наклон $\pm 45^{\circ}$.



Jum.: [1] Pe n rose R., "Proc. Roy. Soc. London, Ser. A», 1965, v. 284, № 1397, p. 159-203; [2] Kruskal M., "Phys. Rev.», 1960, v. 119, № 5, p. 1743-45. A.Т. Ильичев. ПEРВАЯ КВАДРАТИЧНАЯ ФОРМА п о в е р х н о с т и $M^{k}$ в $n$-мерном евклидовом пространстве $\mathbb{R}^{n}$, радиус-вектор текущей точки которой задан параметрически $\boldsymbol{r}=\boldsymbol{r}\left(u^{1}, \ldots\right.$, $\left.u^{k}\right),-$ форма вида



то есть



\[

(d \boldsymbol{r}, d \boldsymbol{r})=g_{a b} d u^{a} d u^{b}

\]



\[

g_{a b}=\left(\frac{\partial \boldsymbol{r}}{\partial u^{a}}, \frac{\partial \boldsymbol{r}}{\partial u^{b}}\right) \text {. }

\]



В инвариантной форме П. к. Ф. в точке $m \in M^{k}-$ это ограничение евклидовой метрики $d s^{2}=d x_{1}^{2}+\ldots+d x_{n}^{2}$ на касательное пространство $T_{m} M^{k}$. В случае когда $n=3$, часто используют классич. обозначения



\[

E=g_{11}, F=g_{12}, G=g_{22} \text {. }

\]



Коэффициенты $g_{i j}$ образуют тензор, называемый м е т рическим тензором.



См. также Риманова геометрия.



В. В. Трофимов.



ПЕРВАЯ КРАЕВАЯ ЗАДАЧА - см. Краевая задача, Дирихле задача.



ПЕРВИЧНАЯ СВЯЗЬ - см. Гамильтонов формализм.



ПЕРВОЕ НАЧАЛО ТЕРМОДИНАМИКИ - см. Термодинамики законы.



IEPBOE УРАВНЕНИЕ КОЛМОГОРОВА - см. Колмогорова уравнение.





\end{document}